\section{Introduction}\label{sec:introduction}

For a smooth connected projective complex surface \(S\), write
\[
p_g(S)=h^0(S,\omega_S),\qquad q(S)=h^1(S,\mathcal O_S).
\]
The Chow group $\CH_0(S)$ is the group of integral zero-cycles modulo
rational equivalence, where the relations are generated by divisors
of rational functions on curves in $S$.  The degree of a complex point is one, and
$\CH_0(S)^0$ denotes the kernel of the degree map.
The Albanese homomorphism
\(\operatorname{alb}_S:\CH_0(S)^0\to\Alb(S)(\mathbb C)\) is
surjective. Bloch's conjecture predicts that it is an isomorphism when
\(p_g(S)=0\). We prove the conjecture. The main theorem treats the
case \(q=0\), when the Albanese variety is a point and the prediction
is that any two complex points are rationally equivalent.

\begin{theorem}\label{thm:bloch}
Let \(S\) be a smooth connected projective surface over \(\mathbb C\)
with \(p_g(S)=q(S)=0\). Then
\[
\deg:\CH_0(S)\longrightarrow\mathbb Z
\]
is an isomorphism, where \(\CH_0(S)\) has integral coefficients.
\end{theorem}

Theorem~\ref{thm:bloch} imposes no minimality, fundamental-group, or
surface-construction hypothesis. Combining it with the classical
non-general-type case gives the full statement.

\begin{corollary}[Bloch's conjecture for surfaces with \(p_g=0\)]\label{cor:bloch-pg-zero}
Let \(S\) be a smooth connected projective surface over \(\mathbb C\)
with \(p_g(S)=0\). Then the Albanese homomorphism
\[
\operatorname{alb}_S:\CH_0(S)^0\longrightarrow\Alb(S)(\mathbb C)
\]
is an isomorphism of abelian groups, where \(\CH_0(S)\) has integral
coefficients.
\end{corollary}

For a surface of general type, $p_g=0$ forces $q=0$, so
Theorem~\ref{thm:bloch} applies. The other cases are the theorem of
Bloch--Kas--Lieberman \cite{BKL}. We give the precise reduction in
Section~\ref{subsec:integral-proof}.


These zero-cycle statements also have consequences for rational Chow
motives. For $p_g=q=0$, every rational Chow class on a power of $S$
or a Hilbert scheme of points is determined by its cohomology class.
We state the precise consequences and their motivic antecedents after
the proof overview.

\subsection*{History and antecedents}

Mumford showed that a nonzero holomorphic two-form prevents the
zero-cycles on a surface from admitting a finite-dimensional
parametrization modulo rational equivalence \cite{Mumford}.
Bloch's conjecture predicts the converse at the level of the Albanese
kernel: if $p_g=0$, degree and the Albanese image should account for
all rational-equivalence classes of zero-cycles \cite{BlochLectures}.
Bloch, Kas, and Lieberman proved this integral assertion for every
surface of Kodaira dimension less than two \cite[pp.~135--136]{BKL}.
Thus the remaining general-type problem concerns $p_g=q=0$, where
the Albanese variety carries no information and the degree map is
predicted to determine a zero-cycle completely.

Several successful approaches use additional geometry of the surface.
The group-action method of Inose--Mizukami \cite{InoseMizukami} and
its extension by Barlow \cite{Barlow} compare zero-cycles with their
images on suitable quotient surfaces. Bauer's treatment of Inoue
surfaces with $K^2=7$ gives a later application, expressed through
relations among the quotient correspondences
\cite[Theorem~0.2 and Section~1]{BauerInoue}.
Pedrini--Weibel recast involution arguments using the transcendental
motive: for a bidouble cover, vanishing of the transcendental motives
of the three intermediate quotients forces its own vanishing
\cite[Theorem~2.5]{PedriniWeibel}. Their applications include
numerical Godeaux and classical Campedelli examples, Keum--Naie,
Burniat, and Inoue surfaces, giving examples for each $1\le K^2\le8$.

A different route uses families of surfaces and the actions of
correspondences on zero-cycles. For a family of regular surfaces
whose fibered self-product has a rationally connected smooth projective
compactification, Voisin shows that a relative correspondence acting
trivially on holomorphic two-forms acts nilpotently on degree-zero
zero-cycles. She applies this theorem to Catanese surfaces
\cite[Theorems~0.3 and~0.6]{Voisin2014}. Passing to suitable quotients
then gives Bloch's conjecture for Barlow surfaces. This extends the
earlier result for Barlow surfaces carrying extra symmetries to the
general members of that family. The present argument requires no
prescribed family, covering construction, or automorphism of the
surface.

Recent preprints propose proofs with overlapping scope.
Guletski\u\i's Theorem~A assumes algebraicity of total cohomology
and claims vanishing of the Albanese kernel over algebraically closed
fields of characteristic zero \cite[Theorem~A and Section~2.1]{Guletskii}.
Over $\mathbb C$, this cohomological assumption is equivalent to
$p_g=q=0$, by Hodge theory and the Lefschetz $(1,1)$ theorem, so it
covers the full scope of Theorem~\ref{thm:bloch}. His proposed argument studies a hypersurface section through the
diagonal in the square of the surface, together with a nilpotence
argument for correspondences.
Banerjee proposes a separate proof for minimal general-type surfaces
with $p_g=q=0$ and $K^2=9$, using pluricanonical curves and their
Jacobians \cite[Theorems~1.1 and~2.1]{BanerjeeRankOne}.
We cite these as proposed proofs; the argument developed here obtains
its Chow relation from virtual localization on a Quot scheme.

The use of the diagonal to control zero-cycles has its antecedents in
Bloch--Srinivas \cite{BlochSrinivas}. Surface Quot obstruction theories
and localization relations were developed by Marian--Oprea--Pandharipande
\cite{MOP}, with subsequent virtual intersection theory in \cite{OP}
and its $K$-theoretic counterpart in \cite{AJLOP}.
Ciocan-Fontanine--Kapranov constructed derived Quot schemes \cite{CFK}.
We use the derived moduli and determinant constructions of
To\"en--Vaqui\'e and Sch\"urg--To\"en--Vezzosi \cite{TV,STV}
to impose the determinant constraint. The Chow calculation adapts the
nested-Hilbert-scheme recursion of Ellingsrud--G\"ottsche--Lehn
\cite{EGL} by leaving two surface factors unintegrated throughout.
This is the step that produces a correspondence identity rather than
only a numerical intersection formula.

\subsection*{The mechanism of the proof}

The rational cohomology of a surface with $p_g=q=0$ is generated by algebraic
classes, but cohomological identities alone do not settle rational
equivalence. After the classical and birational reductions, we work on
a minimal surface of general type or a point blowup $X$ of that surface.
The objective is a relation
\begin{equation}\label{eq:intro-diagonal}
 0=c[\Delta_X]+\Gamma\quad\text{in }\CH^2(X\times X)_{\mathbb Q},
\end{equation}
where $c\ne0$ and $\Gamma$ is a sum of external products of cycles.
The diagonal acts as the identity on zero-cycles, whereas every external
product of total codimension two kills degree-zero zero-cycles.
Thus the relation gives $\CH_0(X)^0\otimes\mathbb Q=0$.
Constructing this Chow relation and detecting its coefficient are the
two main stages.

For line bundles $L_1,L_2,D$, consider rank-zero quotients of
$L_1\oplus L_2$ whose determinant has Picard class $D$.
Opposite torus weights on the summands give a localization formula.
We construct its obstruction theory and the compatible fixed theory
using a derived determinant fiber normalized at a point of $X$
(Sections~\ref{sec:derived} and~\ref{sec:fixed}). A fixed quotient
splits into a quotient of each line bundle, each described by an effective
divisor and a zero-dimensional subscheme. The divisor line bundles
$T,H$ satisfy $T+H=D$; on an irregular surface this determinant
condition also relates their Picard parameters.

Two second-Chern-character insertions of the universal quotient leave
a class on $X\times X$. We choose the virtual dimension so that
equivariant degree forces its codimension-two part to vanish. To calculate it in Chow, we retain both output copies
of $X$ throughout the Hilbert-scheme recursion of
Ellingsrud--G\"ottsche--Lehn~\cite{EGL}. The recursion replaces the
point incidence classes by diagonals identifying surface coordinates.
Terms joining both outputs contribute a multiple of $[\Delta_X]$;
terms with separated outputs contribute external products. For each
fixed fiber expression, this procedure chooses its diagonal coefficient
from the intersection numbers of its line bundles and the Chern
numbers of the surface. The total $c$ is the weighted sum over the
fixed contributions that occur (Section~\ref{sec:relation}). The rule
does not require a diagonal decomposition on an individual surface
to be unique.

For a minimal general-type surface with $p_g=q=0$ and canonical
class $K$, one has $b_2=10-K^2$. When $K^2\le8$, its lattice has a direction orthogonal
to the canonical class. A lattice construction and at most three
point blowups provide a class $\alpha$ with $\alpha^2<0$ and Quot
data with the following property. For each contributing fixed pair
with $\alpha H=0$, pairing with $\alpha\boxtimes\alpha$ kills
its external terms. Every remaining contributing pair has zero
point lengths, so its chosen
diagonal coefficient is zero, while its pairing is $(\alpha H)^2$.
Consequently the full relation gives
\[
 0=c\alpha^2+\sum_{\alpha H\ne0}(\alpha H)^2.
\]
The sum runs over a finite set of effective fixed pairs, and a separate
effectivity argument shows that at least one occurs. Since
$\alpha^2<0$, this proves $c>0$ (Section~\ref{sec:nonrank}).

When $K^2=9$, the minimal surface has no such orthogonal direction.
We choose $D=3K$ and arrange that only the numerical divisor pairs
$(3K,0)$ and $(2K,K)$ can contribute to the diagonal coefficient.
Write $c_{\mathrm{out}}$ and $c_{\mathrm{mid}}$ for their respective
contributions, including both orders and all admissible point lengths.
They can be compared on the Cartwright--Steger surface, which has the same
Chern numbers but $p_g=q=1$. The real part of its holomorphic two-form
has nonzero square and is orthogonal to all divisor classes, so pairing
detects the diagonal coefficient. The divisor cosection of
Chang--Kiem~\cite{ChangKiem}, with Kiem--Li
localization~\cite{KiemLi}, reduces its fixed contributions to fibers
over the empty and canonical divisors. These have virtual
multiplicities $1$ and $-s$ with $s>0$. After restricting to those
fibers, the universal rule compares the same two coefficients on both
surfaces and gives $c_{\mathrm{out}}=s c_{\mathrm{mid}}$.
A length-two Hilbert-scheme calculation gives
$c_{\mathrm{mid}}=12$. On the original surface the two coefficients
occur with nonnegative multiplicities, and the outer type occurs once,
so its total coefficient is positive (Section~\ref{sec:rankone}).

Finally, rational vanishing makes each integral degree-zero class
torsion. Roitman's torsion theorem, together with $q=0$, makes it
zero~\cite{BlochRoitman,Roitman}.

The technical constructions used to obtain
\eqref{eq:intro-diagonal} are recorded separately. They include
compatible obstruction theories for the based determinant constraint
(Propositions~\ref{prop:quot-pot} and~\ref{prop:fixed-pot}),
an excess comparison for rank-one Quot schemes
(Lemma~\ref{fix:excess}), and a Chow-valued
universality statement with two marked surface factors. The latter
applies to the tautological pushforwards specified in
Proposition~\ref{prop:universal}, independently of the numerical choices
made for the zero-cycle application.

Figure~\ref{fig:proof-flow} records how the numerical branches feed
into the same diagonal criterion.

\begin{figure}[ht]
\centering
\small
\begin{tabular}{c}
\fbox{\begin{minipage}{0.86\linewidth}\centering
Two marked surface factors in a determinant-fixed Quot calculation\\
Sections~\ref{sec:derived}--\ref{sec:relation}
\end{minipage}}\\[3pt]
$\Downarrow$\\[3pt]
$0=c[\Delta_X]+\Gamma$ in $\CH^2(X\times X)_{\mathbb Q}$,
with $\Gamma$ a sum of external products\\[8pt]
\begin{tabular}{c@{\hspace{1.4em}}c}
\fbox{\begin{minipage}{0.39\linewidth}\centering
$K_{X_0}^2\leq 8$\\
Lattice class and at most three blowups\\
$\alpha^2<0$,\quad $0=c\alpha^2+\sum(\alpha H)^2$\\
A nonempty sum forces $c>0$\\
Section~\ref{sec:nonrank}
\end{minipage}}
&
\fbox{\begin{minipage}{0.39\linewidth}\centering
$K_{X_0}^2=9$\\
Cartwright--Steger comparison\\
$c_{\mathrm{out}}=s c_{\mathrm{mid}}$,\quad $c_{\mathrm{mid}}=12$\\
$c=12(s+\tau_{X_0})>0$\\
Section~\ref{sec:rankone}
\end{minipage}}
\end{tabular}\\[3pt]
$\Downarrow$\\[3pt]
$\CH_0(X_0)^0\otimes\mathbb Q=0$\\[3pt]
$\Downarrow\quad\text{Roitman's torsion theorem and }q=0$\\[3pt]
$\deg:\CH_0(S)\xrightarrow{\sim}\mathbb Z$
\end{tabular}
\caption{The general-type proof, after passage to the minimal model
$X_0$ of $S$.  Each branch chooses Quot data and proves that the
resulting diagonal coefficient is nonzero. In the left branch,
the displayed squares come from the external-product terms with
$\alpha H\ne0$; those terms have zero diagonal coefficient.
In the right branch, the outer and middle coefficients are defined
by the same numerical rule on the target and auxiliary surfaces;
$-s$ is the canonical-divisor virtual coefficient on the auxiliary
surface, with $s>0$, and $\tau_{X_0}$ counts the nontrivial torsion
line bundles.  Birational invariance returns the conclusion to $S$.}
\label{fig:proof-flow}
\end{figure}


\subsection*{Motivic consequences}

The full zero-cycle statement also has consequences for rational Chow
motives. Write \(h(Y)_{\mathbb Q}\) for the rational Chow motive of a
smooth projective complex variety \(Y\), and \(t_2(S)\) for the
transcendental degree-two summand defined by Kahn--Murre--Pedrini.
Let \(\mathbb L\) denote the effective Lefschetz motive, whose Betti
realization is \(\mathbb Q(-1)\) in degree two. A rational Chow motive
is \emph{pure Tate} if it is a finite direct sum of powers of
\(\mathbb L\).

\begin{corollary}[Motivic consequences]\label{cor:tate-motives}
Let \(S\) be a smooth connected projective surface over \(\mathbb C\)
with \(p_g(S)=0\). Then \(h(S)_{\mathbb Q}\) is finite-dimensional in
the sense of Kimura and \(t_2(S)=0\).

If, in addition, \(q(S)=0\), then
\[
h(S)_{\mathbb Q}\simeq
\mathbf 1\oplus\mathbb L^{\oplus b_2(S)}\oplus\mathbb L^2.
\]
In this case, for every \(m,n\ge0\), the rational Chow motives of the
power \(S^m\) and the Hilbert scheme \(S^{[n]}\) of length-\(n\)
zero-dimensional subschemes are pure Tate, with
\(S^0=S^{[0]}=\operatorname{Spec}\mathbb C\).
For \(Y=S^m\) or \(Y=S^{[n]}\), the rational cycle-class maps satisfy
\[
\operatorname{cl}_Y^i:\CH^i(Y)_{\mathbb Q}
 \xrightarrow{\;\sim\;}H^{2i}(Y,\mathbb Q),
\qquad H^{2i+1}(Y,\mathbb Q)=0
\quad\text{for every }i\ge0.
\]
\end{corollary}

These are consequences of established motivic criteria. In the framework
of Kimura's finite-dimensional motives \cite{Kimura},
Guletski\u\i--Pedrini proved the equivalence with Bloch's conjecture
for surfaces with $p_g=0$ and recorded the resulting Tate decomposition when
$q=0$ \cite[Theorem~27 and Remark~28]{GuletskiiPedrini}.
The surface decomposition of Kahn--Murre--Pedrini identifies vanishing
of the Albanese kernel with $t_2(S)=0$ \cite{KMP}; the Hilbert-scheme
assertion then uses
the motive formula of de Cataldo--Migliorini \cite{deCataldoMigliorini}.
In particular, every rational Chow class on the stated powers and Hilbert
schemes is determined by its cohomology class. This conclusion uses
the Tate decomposition, not finite-dimensionality alone.
Section~\ref{subsec:motivic-consequences} gives the proof.

Section~\ref{sec:reductions}
first passes to minimal surfaces of general type.
Section~\ref{sec:derived} constructs the determinant-fixed theory
and its fixed part; Section~\ref{sec:fixed} computes the resulting
virtual classes and the divisor cosection.  Section~\ref{sec:relation}
proves the marked Chow relation and its action on zero-cycles.
Sections~\ref{sec:nonrank} and~\ref{sec:rankone} prove the two
nonvanishing statements. Section~\ref{subsec:integral-proof} gives the
integral conclusion, and Section~\ref{subsec:motivic-consequences}
derives the motivic consequences.

\subsection*{Conventions}

All schemes and cohomology groups are over \(\mathbb C\).
The definitions of \(\CH_0(S)\) and \(\operatorname{alb}_S\) above,
Theorem~\ref{thm:bloch}, Corollary~\ref{cor:bloch-pg-zero} and its proof,
Lemma~\ref{lem:birational}, and Section~\ref{subsec:integral-proof} use
integral Chow groups. Elsewhere Chow groups have rational coefficients;
in particular, all Chow motives in Section~\ref{subsec:motivic-consequences}
have rational coefficients.
We use cohomological indexing for complexes. Line bundles are written
additively; in intersections the same symbol denotes their first
Chern class. For a perfect complex, its Euler class with a nonzero
torus weight is interpreted in localized equivariant Chow theory.
Statements about virtual classes always use the specified perfect
obstruction theories. A comparison of their classes in \(K\)-theory
is used to compute Chern classes only after the relevant obstruction
bundle or compatible map of theories has been identified.
